\section{Artifact Appendix}
\label{app:artifact}

\subsection{Abstract}

The artifact provides Accelgen, the compiler framework for \name{}. Starting
from PyTorch model code, the workflow generates Linalg-based MLIR, optimizes
Linalg operations, schedules kernels, and lowers the scheduled computation to
a Spatial Architecture Primitive (SAP) graph for performance modeling. The
outputs include intermediate MLIR and performance metrics such as latency,
energy, FLOPs, throughput, and energy efficiency. These outputs can be
inspected independently of the figures and tables in the paper.

\subsection{Artifact Checklist}

The artifact produces the following outputs to support reproducibility and
inspection of each compilation stage.
\begin{itemize}[leftmargin=2em,itemsep=0.25em]
  \item \textbf{Generated Linalg MLIR.} The Linalg representation generated
    from the input model.
  \item \textbf{Optimized Linalg MLIR.} The Linalg representation after
    Linalg-level optimizations.
  \item \textbf{Scheduled Clusters.} Computation clusters with their schedules
    and associated dataflow parameters.
  \item \textbf{SAP Graph.} The generated SAP graph representing the
    accelerator architecture.
  \item \textbf{Performance Report.} Modeled performance metrics for the
    generated accelerator.
\end{itemize}

\subsection{Installation}

We provide two Docker-based methods for installing and running \name{}.

\textbf{Using the Prebuilt Docker Image (Recommended).} Pull the prebuilt
image and launch a container with the repository mounted as the working
directory.

\begin{verbatim}
$ git clone https://github.com/XiaoyeAsZ/accelgen.git
$ cd accelgen
$ docker pull xiaoyeasz/accelgen:latest
$ docker run --gpus all -it --rm \
    -v $(pwd):/workspace/accelgen \
    -w /workspace/accelgen \
    xiaoyeasz/accelgen:latest
\end{verbatim}

\textbf{Building from the Dockerfile.} Alternatively, build the image locally
and launch a container using the same repository mount.

\begin{verbatim}
$ git clone https://github.com/XiaoyeAsZ/accelgen.git
$ cd accelgen
$ docker build -f accelgen.dockerfile -t accelgen:latest .
$ docker run --gpus all -it --rm \
    -v $(pwd):/workspace/accelgen \
    -w /workspace/accelgen \
    accelgen:latest
\end{verbatim}

\subsection{Experimental Workflow}

The workflow consists of five stages that can be run separately to inspect
and reuse intermediate outputs. Each stage requires the output of the
preceding stage. Two evaluation modes are available: \texttt{rapid} for a
quick evaluation of the block~0 FFN workload in Llama3-8B, and \texttt{full}
for the complete evaluation suite. The full mode can take substantially
longer, particularly during kernel scheduling.

Inside the container, enter the artifact evaluation directory:

\begin{verbatim}
$ cd ae
\end{verbatim}

In the commands below, replace \texttt{MODE} with \texttt{rapid} or
\texttt{full}. All output directories are relative to
\texttt{ae/results/MODE/} in the repository.

\begin{enumerate}[leftmargin=2.4em,itemsep=0.45em]
  \item \textbf{Generate MLIR.}\\
  {\small \texttt{\$ bash MODE/01\_generate\_mlir.sh}}\\
  The script invokes \texttt{benchmark/dump\_mlir.py} for each selected
  combination of model, action, block, layer, batch size, and sequence
  length. The generated MLIR files are written to \texttt{mlir/}.

  \item \textbf{Optimize Linalg.}\\
  {\small \texttt{\$ bash MODE/02\_linalg\_opt.sh}}\\
  The optimizer performs mixed-precision resolution, named-operation
  generalization, generic marking, dead-operation elimination, generic
  fusion, and tensor-operation folding. The optimized MLIR files are
  written to \texttt{generic/}.

  \item \textbf{Schedule Kernels.}\\
  {\small \texttt{\$ bash MODE/03\_kernel\_schedule.sh}}\\
  The \texttt{func.func(kernel-schedule\{...\})} pass searches over tiling,
  loop, unrolling, and cluster choices using the selected edge or server
  resource configuration. The scheduled MLIR files are written to
  \texttt{scheduled/}. By default, each cluster contains at most six
  operations (\texttt{MAX\_OPERATION=6}).

  \item \textbf{Lower to DAP.}\\
  {\small \texttt{\$ bash MODE/04\_lower\_to\_dap.sh}}\\
  The \texttt{convert-to-dap\{config-path=...\}} pass lowers each scheduled
  workload to the hardware-facing DAP dialect. The resulting MLIR files
  are written to \texttt{dap/} and contain operations such as
  \texttt{dap.sram}, \texttt{dap.data\_path}, arithmetic operations, and
  routing operations.

  \item \textbf{Model Performance.}\\
  {\small \texttt{\$ bash MODE/05\_model\_performance.sh}}\\
  The \texttt{model-performance\{config-path=...\}} pass estimates latency,
  energy, and FLOPs and emits annotated MLIR files in \texttt{performance/}.
  The accompanying \texttt{logs/performance.log} contains workload headers
  and metric records suitable for processing with the repository's
  analysis scripts.
\end{enumerate}
